% Created 2026-10-01 Thu 17:18
% Intended LaTeX compiler: pdflatex
\documentclass[11pt]{article}

\usepackage[utf8]{inputenc}
\usepackage[T1]{fontenc}
\usepackage{graphicx}
\usepackage{longtable}
\usepackage{wrapfig}
\usepackage{rotating}
\usepackage[normalem]{ulem}
\usepackage{amsmath}
\usepackage{amssymb}
\usepackage{capt-of}
\usepackage{hyperref}
\author{Colton Quirk}
\date{\today}
\title{How to Simulate a Galaxy}
\hypersetup{
 pdfauthor={Colton Quirk},
 pdftitle={How to Simulate a Galaxy},
 pdfkeywords={},
 pdfsubject={},
 pdfcreator={},
 pdflang={English}}
\begin{document}

\maketitle
\tableofcontents

\section{Introduction}
\label{sec:org282501a}
This is my exploration into simulatin a galaxy using Julia.
My goal is to build up a simulation using more and more advanced features as I add to it.
This will allow me to learn about the process of simulating a galaxy, as well as best practices for high performance Julia code.

I will start with:
\begin{itemize}
\item Basic particle type
\item Direct summation of forces
\item Runge-Kutta and Symplectic integrators with fixed and/or variable timesteps
\end{itemize}

Eventually I will add things like:
\begin{itemize}
\item Barnes-Hut algorithm for force summation
\item Fast Multipole Method for force summation
\item Smoothed Particle Hydrodynamics
\item Mesh-grid fluid-like simulations
\end{itemize}

Finally/during the development I will learn how to make these systems parallel.
This could include, but is not limited to, multi-threading and GPU acceleration.
\subsection{Setting up the package}
\label{sec:org90dd33a}

\begin{verbatim}
julia --eval 'using Pkg; Pkg.generate("GalaxySim")'

cd GalaxySim
julia --eval 'using Pkg; Pkg.activate("."); Pkg.instantiate()'

cd GalaxySim
touch main.jl
echo using GalaxySim >> main.jl

git init .
git add Project.toml
git add main.jl
git add src/
git commit -m 'Initial Commit'
\end{verbatim}
\subsection{{\bfseries\sffamily TODO} Configure startup.jl}
\label{sec:org2c84c71}
\textbf{Remember to activate the environment!}
\begin{itemize}
\item Add this to \texttt{startup.jl}
\end{itemize}
\section{Particle Types}
\label{sec:org7064015}
\subsection{Basic Particle}
\label{sec:orgfbbd58a}
A basic particle needs to hold a few key pieces of information.
\begin{enumerate}
\item Its mass, to determine how it effects other particles
\item Its position, so we know where it is relative to the other particles
\item Its velocity, so we know how it moves throughout the galaxy.
\end{enumerate}

I in previous projects have been very fond of the Dynamic Form of ODEs as described by \cite{landauComputationalPhysics2024}.

\begin{verbatim}
using GalaxySim
\end{verbatim}
\subsection{Smoothed Particle Hydrodynamics Particle}
\label{sec:orge14667e}
\section{Summation of Forces}
\label{sec:orgad818a3}
\subsection{Direct Summation}
\label{sec:org634afff}
\subsection{Barnes-Hut Algorithm}
\label{sec:orgb8ecda3}
\subsection{Fast Multipole Method}
\label{sec:orga89d195}
\section{Integration of Forces}
\label{sec:org7bf01f9}
\subsection{Euler Method}
\label{sec:org1da0f17}
\subsection{Symplectic Methods}
\label{sec:org27b005c}
\section{Fluid Dynamics}
\label{sec:orgc96882d}
\subsection{Smoothed Particle Hydrodynamics}
\label{sec:orga2f2b1e}
\[
\rho(\mathbf{r}) = \sum_{b=1}^{N_{\mathrm{neigh}}} m_{b} W(\mathbf{r}-\mathbf{r}_b, h_b)
\]
\subsection{Grid based methods}
\label{sec:org8897bef}

\bibliographystyle{mnras}
\bibliography{references}
\end{document}
